\subsection{单模的类}

用 \(\mathrm{Is}_{\mathrm{A}}\{\mathrm{X},\mathrm{Y}\}\) 表示关系

``A 是环且 X、Y 是同构的 A-模。''

这是关于 X 与 Y 的一个等价关系。若 X 是一个 A-模，则我们把对于 \(\operatorname{Is}_{\mathrm{A}}\) 与 X 等价的对象组成的类记作 \(\operatorname{cl}(\mathrm{X})\)，并称之为 A-模 X 的类（《集合论》，II, §6, No.~9, 第 122 页）。按定义，\(\operatorname{cl}(\mathrm{X})\) 是一个同构于 X 的 A-模；此外，两个 A-模 X 与 Y 同构当且仅当我们有 \(\operatorname{cl}(\mathrm{X}) = \operatorname{cl}(\mathrm{Y})\) 。

设 A 是环。关系

``\(\lambda\) 是一个有限生成 A-模的类''

在 \(\lambda\) 中是集合化的（《集合论》，II, §1, No.~4, 第 68 页）。事实上，每个有限生成 A-模都同构于形如 \(A_{s}^{n}/R\) 的 A-模，其中 n 是自然数，R 是 \(A_{s}^{n}\) 的一个子模，于是我们的断言由《集合论》，II, §6, No.~9, 第 122 页得出。

我们用 \(\mathcal{F}(A)\) 表示有限生成 A-模的类组成的集合。每个单 A-模都是单生成的（VIII，第 46 页，命题 1），从而单 A-模的类构成 \(\mathcal{F}(A)\) 的一个子集，从现在起我们把它记作 \(\mathcal{S}(A)\)（或简作 S）。当环 A 交换时，映射 \(\mathfrak{m} \mapsto \operatorname{cl}(A/\mathfrak{m})\) 是从 A 的极大理想组成的集合到集合 \(\mathcal{S}(A)\) 的一个双射（出处同前及 VIII，第 46 页，注 1）。当 A 是左阿廷环时，集合 \(\mathcal{S}(A)\) 有限（VIII，第 51 页，注 b)）。

